\documentclass[12pt,a4paper,UTF8]{ctexart}
\usepackage[a4paper,top=22mm,bottom=22mm,left=22mm,right=22mm,headheight=15pt]{geometry}
\usepackage{amsmath,amssymb,mathtools,bm}
\usepackage{graphicx,xcolor,booktabs,longtable,array,enumitem,fancyhdr,caption,float,microtype,setspace}
\newcommand{\Needspace}[1]{\par\begingroup\dimen0=#1\relax\dimen2=\pagegoal\advance\dimen2 by-\pagetotal\ifdim\dimen0>\dimen2\newpage\fi\endgroup}
\newcommand{\needspace}[1]{\par\begingroup\skip0=\lastskip\vskip-\skip0\dimen0=#1\relax\advance\dimen0 by\skip0\vskip\dimen0\penalty-100\vskip-\dimen0\vskip\skip0\penalty9999\endgroup}
\usepackage[hidelinks,unicode]{hyperref}
\usepackage[most]{tcolorbox}
\definecolor{ContestBlue}{HTML}{17365D}
\hypersetup{pdftitle={第 050 套},pdfauthor={}}
\setstretch{1.24}
\setlength{\parindent}{2em}
\setlength{\parskip}{0.35em}
\setlength{\emergencystretch}{2em}
\setlist{itemsep=0.35em,topsep=0.4em}
\captionsetup{font=small,labelfont=bf,justification=centering}
\setcounter{secnumdepth}{0}
\pagestyle{fancy}
\fancyhf{}
\fancyhead[L]{\small 第 050 套}
\fancyfoot[C]{\thepage}
\renewcommand{\headrulewidth}{0.4pt}
\ctexset{
  section={format=\Large\bfseries\color{ContestBlue},beforeskip=1.2em,afterskip=0.7em},
  subsection={format=\large\bfseries\color{ContestBlue},beforeskip=1.0em,afterskip=0.5em},
  subsubsection={format=\normalsize\bfseries,beforeskip=0.8em,afterskip=0.35em}
}
\makeatletter
\newcommand{\tightlist}{\setlength{\itemsep}{0pt}\setlength{\parskip}{0pt}}
\newcommand{\maxwidth}{\ifdim\Gin@nat@width>\linewidth\linewidth\else\Gin@nat@width\fi}
\newcommand{\maxheight}{\ifdim\Gin@nat@height>0.72\textheight 0.72\textheight\else\Gin@nat@height\fi}
\newcommand{\pandocbounded}[1]{\begingroup\setkeys{Gin}{width=\maxwidth,height=\maxheight,keepaspectratio}#1\endgroup}
\makeatother
\begin{document}
\begin{titlepage}
  \centering
  \vspace*{58mm}
  {\Huge\bfseries\color{ContestBlue} 第 050 套\par}
  \vfill
  \vspace*{28mm}
\end{titlepage}
\section{第 1 题}\label{ux7b2c-1-ux9898}

空气以速度 \(v_0\) 竖直向下流向挡板，其中比例为 \(\eta=0.1\)
的空气分子被挡板近似原速弹回。把颗粒物先近似为直径和高度均为
\(2r\)、密度为 \(\rho\)
的均匀圆柱，圆柱轴线始终竖直；空气和颗粒在进入挡板附近时均以 \(v_0\)
向下运动。反弹分子与颗粒发生完全弹性碰撞，且只有在颗粒距挡板不超过
\(h_0\)
时碰撞作用才显著。忽略颗粒与其他空气分子、颗粒之间的相互作用。已知 \[
v_0\sim10\,\mathrm{m\,s^{-1}},\qquad h_0=1.0\ {\rm cm},\qquad
\rho_0=1.29\ {\rm kg\,m^{-3}},\qquad
\rho=2.0\times10^3\ {\rm kg\,m^{-3}}.
\]

（1）估算空气分子碰撞使圆柱颗粒获得的加速度数量级；

（2）求能够随气流离开的圆柱颗粒的最大半径 \(r_0\)；

（3）若颗粒改为半径为 \(r\)
的球，空气分子按镜面反射规律与球面碰撞，求最大半径 \(r_0'\)。

\begin{figure}[H]
\centering
\includegraphics[width=0.778\linewidth,height=\textheight,keepaspectratio,alt={图1-1　颗粒与挡板反射气流的相对运动}]{../figures/question-01.pdf}
\caption*{图1-1　颗粒与挡板反射气流的相对运动}
\end{figure}

\newpage

\section{第 2 题}\label{ux7b2c-2-ux9898}

质量为 \(M\)、长度为 \(L\)
的均匀细杆只能在竖直平面内运动，杆的两端始终与半径为 \(R\)
的光滑半圆形碗内壁接触，且 \(R>1.5L\)。

（1）求细杆在最低平衡位置附近的微小振动周期 \(T\)；

（2）若振动的最大偏角为
\(\theta_m\)，碗对细杆的合支持力大小在转折点与通过平衡点时之差写成 \[
\Delta N=N\big|_{\theta=\theta_m}-N\big|_{\theta=0}
=\alpha Mg\theta_m^2,
\] 求 \(\alpha\)。只保留 \(\theta_m\) 的二阶小量。

\begin{figure}[H]
\centering
\includegraphics[width=0.8\linewidth,height=\textheight,keepaspectratio,alt={图2-1　细杆作为圆内定长弦的几何约束}]{../figures/question-02.pdf}
\caption*{图2-1　细杆作为圆内定长弦的几何约束}
\end{figure}

\newpage

\section{第 3 题}\label{ux7b2c-3-ux9898}

同心的内、外圆边界半径分别为 \(R_1\)、\(R_2\)，其中
\(R_1<R_2\)。环形区域中有垂直纸面向里的匀强磁场
\(B_0\)；内圆区域中另有垂直纸面的匀强磁场 \(B_x\)。规定垂直纸面向里为
\(B_x\) 的正方向。质量为 \(m\)、电荷量为 \(q>0\)
的粒子从外圆边界沿半径方向射入，初速度为 \(v_0\)。

（1）求粒子刚能触及内圆边界的临界初速度 \(v_{0\min}\)；

（2）若 \(v_0>v_{0\min}\)，调节 \(B_x\) 使粒子进入内圆后恰好通过圆心，求
\(B_x\)。

\begin{figure}[H]
\centering
\includegraphics[width=0.473\linewidth,height=\textheight,keepaspectratio,alt={图3-1　同心圆磁场区域与径向入射条件}]{../figures/question-03.pdf}
\caption*{图3-1　同心圆磁场区域与径向入射条件}
\end{figure}

\newpage

\section{第 4 题}\label{ux7b2c-4-ux9898}

（1）半径为 \(r\)、质量为 \(m\) 的均匀圆环在水平地面上以质心速度 \(v_1\)
向右纯滚动，撞到高度为 \(h\) 的台阶，其中
\(h<r/2\)。碰撞完全非弹性，求圆环能够翻过台阶并继续向前滚动时 \(v_1\)
应满足的条件。

（2）可变形圆环纯滚动时受到滚动阻力偶矩 \(M=kN\)，其中 \(N\)
为地面对圆环的压力，\(k\) 为已知长度。若初速度为 \(v_0\)，求停止时间。

（3）在摩擦足够大的水平地面上，一根铁棍从左侧推动上述圆环，使其向右做匀速纯滚动。铁棍与圆环间的动摩擦因数为
\(\mu\)，忽略二者间的滚动摩擦。设接触点半径与竖直向下方向夹角为
\(\theta\)。求使铁棍推力最小时的 \(\theta\) 及最小推力。

\begin{figure}[H]
\centering
\includegraphics[width=0.8\linewidth,height=\textheight,keepaspectratio,alt={图4-1　滚环与台阶碰撞的几何关系}]{../figures/question-04.pdf}
\caption*{图4-1　滚环与台阶碰撞的几何关系}
\end{figure}

\newpage

\section{第 5 题}\label{ux7b2c-5-ux9898}

三根竖直试管的横截面积依次为
\(S\)、\(2S\)、\(S\)，底部连通且连通部分体积可忽略。管中装有密度为
\(\rho\) 的理想流体，平衡时液柱高度均为
\(l\)。忽略黏滞，并假定同一根试管内各处流速相同。小扰动后三根液柱高度分别为
\[
l[1+x_1(t)],\qquad l[1+x_2(t)],\qquad l[1+x_3(t)],
\] 其中 \(|x_i|\ll1\)。

（1）写出 \(x_1,x_2,x_3\) 的约束关系；

（2）用 \(x_1,x_3,\dot x_1,\dot x_3\) 表示系统的动能和势能；

（3）初态为 \(x_1=x_2=\varepsilon\)、系统静止，其中
\(\varepsilon\ll1\)，求 \(x_1(t)\)。

\begin{figure}[H]
\centering
\includegraphics[width=0.8\linewidth,height=\textheight,keepaspectratio,alt={图5-1　三根不同截面积连通管的液面位移}]{../figures/question-05.pdf}
\caption*{图5-1　三根不同截面积连通管的液面位移}
\end{figure}

\newpage

\section{第 6 题}\label{ux7b2c-6-ux9898}

温度为 \(T\) 的平衡态理想气体满足麦克斯韦速度分布 \[
dP=\left(\frac{m}{2\pi kT}\right)^{3/2}
\exp\left(-\frac{m v^2}{2kT}\right)
dv_x\,dv_y\,dv_z.
\] 一根水平圆柱形试管长为 \(l\)、横截面积为 \(S\)、质量为
\(M\)，其中密封着分子质量为 \(m\)、温度为 \(T\)、压强为 \(p_0\)
的理想气体。某时刻突然去掉左壁，保留右壁。定义 \[
f(x)=\int_0^x e^{-u^2}\,du.
\]

（1）试管固定，求时刻 \(t\) 前累计与右壁碰撞的分子数；

（2）试管固定，右壁碰撞完全弹性，求时刻 \(t\) 前分子传给右壁的总动量；

（3）试管可自由运动且很重，求最终速度；

（4）把气体换为与温度为 \(T\)
的试管保持热平衡的光子气体，光子与右壁完全弹性碰撞。已知黑体辐射出射度
\(J=\sigma T^4\)，求很重试管的最终速度。

\begin{figure}[H]
\centering
\includegraphics[width=0.8\linewidth,height=\textheight,keepaspectratio,alt={图6-1　左端开启、右壁保留的气体试管}]{../figures/question-06.pdf}
\caption*{图6-1　左端开启、右壁保留的气体试管}
\end{figure}

\newpage

\section{第 7 题}\label{ux7b2c-7-ux9898}

物屏、薄透镜 \(L_1\)、薄透镜 \(L_2\) 和眼睛依次共轴排列。物屏到 \(L_1\)
的距离为 \(l_1\)，两透镜间距为 \(l_2\)。已知 \[
2f_1=l_2,\qquad 2l_2=l_1,\qquad
D_1=3D_2,\qquad f_1=1.5f_2,
\] 且 \(D_1,D_2\ll f_1,f_2\)。

（1）判断哪个透镜是限制入眼光能的孔径光阑；

（2）求眼睛能看到的物屏最大半径 \(r\)；

（3）若物屏上某点发出的、通过第一透镜中心的光线还能进入眼睛，称该点``明亮''，求明亮区域半径
\(R\)。

\begin{figure}[H]
\centering
\includegraphics[width=0.82\linewidth,height=\textheight,keepaspectratio,alt={图7-1　物屏、双透镜与观察眼的轴向布局}]{../figures/question-07.pdf}
\caption*{图7-1　物屏、双透镜与观察眼的轴向布局}
\end{figure}

\newpage

\section{第 8 题}\label{ux7b2c-8-ux9898}

在地面惯性系 \(S\) 中，\(t=0\) 时火车左端 \(L\) 位于 \(x=0\)，右端 \(R\)
位于 \(x=l\)。两端从静止开始，始终具有相同的地面系坐标加速度
\(a\)，均向右运动，且只讨论 \(t<c/a\)。

（1）在地面时刻 \(t\)，求左端瞬时静止系中测得的瞬时加速度 \(a'\)；

（2）以左端本征时 \(\tau=0\) 对应 \(t=0,x=0\)，令
\(\sin\phi=at/c\)，写出左端世界线 \(c\tau(\phi)\) 和 \(x_L(\phi)\)；

（3）以右端本征时 \(\tau'=0\) 对应 \(t=0,x=l\)，写出右端世界线
\(c\tau'(\phi)\) 和 \(x_R(\phi)\)；

（4）地面时刻 \(t_A\) 的左端事件记为 \(A\)。在 \(A\)
对应的瞬时静止系中，求该时刻火车长度，并说明这种运动是否满足刚性运动所要求的相对论性长度关系。

\begin{figure}[H]
\centering
\includegraphics[width=0.8\linewidth,height=\textheight,keepaspectratio,alt={图8-1　两端具有相同坐标加速度的火车世界线}]{../figures/question-08.pdf}
\caption*{图8-1　两端具有相同坐标加速度的火车世界线}
\end{figure}
\end{document}
